\subsection{例}

\begin{enumerate}
\def\labelenumi{\alph{enumi})}
\tightlist
\item
  设 X 是一个拓扑空间。给 \(\mathcal{R}(\mathrm{X})\)，即 X 上（有限）数值函数所成的向量空间，赋予紧收敛拓扑（GT, X, §1, No.~3）：这是使诸限制映射 \(\mathcal{R}(\mathrm{X}) \to \mathcal{R}(\mathrm{H})\) 连续的最粗拓扑（其中 H 取遍 X 的紧子集族，而 \(\mathcal{R}(\mathrm{H})\) 赋予一致收敛拓扑）。III, 第 4 页 的推论 3 表明，\(\mathcal{R}(\mathrm{X})\) 的子集 A 有界，当且仅当对 X 的每个紧子集 H，属于 A 的函数在 H 上的限制所成之集是一致有界的。
\end{enumerate}

* b)（无穷可微函数空间。）设 \(n \geqslant 1\) 是一个整数。对 \(\mathbf{R}^n\) 中的每个开集 U，用 \(\mathcal{C}^\infty(\mathrm{U})\) 表示 U 上无穷可微函数所成的向量空间（VAR, R, 2.3）。设 \(f\) 属于 \(\mathcal{C}^\infty(\mathrm{U})\)。对每个多重指标 \(\alpha = (\alpha_1, \dots, \alpha_n)\)（属于 \(\mathbf{N}^n\)），用 \(\partial^\alpha f\) 表示偏导数 \(\partial^{|\alpha|} f / \partial x_1^{\alpha_1} \dots \partial x_n^{\alpha_1}\)；它是 U 上的一个连续函数（VAR, R, 2.3 与 2.4）。对每个整数 \(m \geqslant 0\) 以及 U 的每个紧子集 H，令

\[
p_{m,\mathrm{H}}(f) = \sup_{\substack{|\alpha |\leqslant m\\ x\in \mathrm{H}}}\left|\partial^{\alpha}f(x)\right|.\tag{1}
\]

于是 \(p_{m,\mathrm{H}}\) 是 \(\mathcal{C}^{\infty}(\mathrm{U})\) 上的一个半范数。

给 \(\mathcal{C}^{\infty}(\mathrm{U})\) 赋予由诸半范数 \(p_{m,\mathrm{H}}\) 所定义的拓扑。这是使诸映射 \(f\to \partial^{\alpha}f\)（从 \(\mathcal{C}^{\infty}(\mathrm{U})\) 到 \(\mathcal{R}(\mathrm{U})\) 中，其中 \(\mathcal{R}(\mathrm{U})\) 赋予紧收敛拓扑）连续的最粗拓扑。存在 U 的紧子集的一个递增序列 \((\mathrm{H}_n)_{n\geqslant 0}\)，它们的内部覆盖 U；半范数族 \(p_{m,\mathrm{H}_n}\) 定义了 \(\mathcal{C}^{\infty}(\mathrm{U})\) 的拓扑，后者于是成为一个可度量化的局部凸空间。空间 \(\mathcal{C}^{\infty}(\mathrm{U})\) 是完备的；换言之，它是一个 Fréchet 空间（II, 第 24 页）：事实上，设 \((f_k)\) 是 \(\mathcal{C}^{\infty}(\mathrm{U})\) 中的一个 Cauchy 序列；对每个 \(\alpha \in \mathbf{N}^n\)，序列 \((\partial^{\alpha}f_k)\) 在完备空间 \(\mathcal{R}(\mathrm{U})\) 中收敛（TG, X, §1, No.~5, 定理 1）于一个连续函数 \(g_{\alpha}\)。对 \(|\alpha|\) 作归纳，我们由 FVR, II, 第 2 页 的定理 1 推出 \(g_{\alpha} = \partial^{\alpha}g_{0}\) 对每个 \(\alpha \in \mathbf{N}^n\) 成立。换言之，序列 \((f_k)\) 收敛于 \(g_0\)（在 \(\mathcal{C}^{\infty}(\mathrm{U})\) 中）。

设 A 是 \(\mathcal{C}^{\infty}(\mathrm{U})\) 的一个子集。要使 A 有界，其充分必要条件是：数 \(\sup_{f\in \mathbf{A}}p_{m,\mathrm{H}}(f)\) 对每个整数 \(m\geqslant 0\) 以及每个

  U 的紧子集 H 而言都是有限的；这一条件意味着：对每个 \(\alpha\in \mathbf{N}^{n}\)，诸函数 \(\partial^{\alpha}f|H\)（\(f\in A\)）所成之集对每个紧的 \(H\subset U\) 都是一致有界的。

设 \(\mathrm{H} \subset \mathrm{U}\) 是紧的。用 \(\mathcal{C}_{\mathrm{H}}^{\infty}(\mathrm{U})\) 表示 \(\mathcal{C}^{\infty}(\mathrm{U})\) 中由支集含于 \(\mathrm{H}\) 的那些函数所组成的子空间。空间 \(\mathcal{C}_{\circ}^{\infty}(\mathrm{U})\)，即 \(\mathrm{U}\) 中具有紧支集的无穷可微函数所成的空间，是诸子空间 \(\mathcal{C}_{\mathrm{H}}^{\infty}(\mathrm{U})\) 的递增有向并，其中 \(\mathrm{H}\) 取遍 \(\mathrm{U}\) 的紧子集族。每个空间 \(\mathcal{C}_{\mathrm{H}}^{\infty}(\mathrm{U})\) 都赋予由 \(\mathcal{C}^{\infty}(\mathrm{U})\) 的拓扑诱导的拓扑，而 \(\mathcal{C}_{\circ}^{\infty}(\mathrm{U})\) 赋予相应的归纳极限拓扑。如果诸集 \(\mathrm{H}_n\) 的内部构成 \(\mathrm{U}\) 的一个覆盖，那么空间 \(\mathcal{C}_{\circ}^{\infty}(\mathrm{U})\) 是诸 Fréchet 空间 \(\mathcal{C}_{\mathrm{H}_n}^\infty (\mathrm{U})\) 的严格归纳极限；因此它是完备的（II, 第 32 页, 命题 9），并且 \(\mathcal{C}_{\circ}^{\infty}(\mathrm{U})\) 的每个有界子集都含于诸子空间 \(\mathcal{C}_{\mathrm{H}_n}^\infty (\mathrm{U})\) 之一（III, 第 5 页, 命题 6）。*

\begin{enumerate}
\def\labelenumi{\alph{enumi})}
\setcounter{enumi}{2}
\tightlist
\item
  （Gevrey 空间。）设 I 是 \(\mathbf{R}\) 中的一个紧区间。对每个整数 \(n \geqslant 0\)，用 \(\mathrm{D}^n f\) 表示定义在 I 上的数值函数的 \(n\) 阶导数（该函数为 \(f\)，只要此导数存在）。设 \(s \geqslant 1\) 与 \(\mathbf{M} \geqslant 0\) 是两个实数。用 \(\mathcal{G}_{s,\mathbf{M}}(\mathrm{I})\) 表示 I 上满足如下条件的那些无穷可微函数 \(f\)（FVR, I, 第 28 页）所成的向量空间：序列 \((|\mathrm{D}^n f| / \mathrm{M}^n (n!)^s)_{n \geqslant 0}\) 在 I 上所有连续函数所成的空间 \(\mathcal{C}(\mathrm{I})\)（赋一致收敛拓扑）中是有界的。空间 \(\mathcal{G}_{s,\mathbf{M}}(\mathrm{I})\) 赋予下述范数后是一个 Banach 空间
\end{enumerate}

\[
\| f \| _ {s, \mathbf {M}} = \sup _ {n \geqslant 0, x \in \mathrm{I}} | \mathrm{D} ^ {n} f (x) | / \mathrm{M} ^ {n} (n!) ^ {s}.
\]

对 \(\mathbf{M} \leqslant \mathbf{M}'\)，我们有 \(\mathcal{G}_{s,\mathbf{M}}(\mathrm{I}) \subset \mathcal{G}_{s,\mathbf{M}'}(\mathrm{I})\)，并且

\[
\| f \| _ {s, \mathbf {M} ^ {\prime}} \leqslant \| f \| _ {s, \mathbf {M}}
\]

对每个 \(f \in \mathcal{G}_{s,\mathsf{M}}(\mathrm{I})\) 成立。用 \(\mathcal{G}_s(\mathrm{I})\) 表示诸空间 \(\mathcal{G}_{s,\mathsf{M}}(\mathrm{I})\) 的并，并给它赋予 \(\mathcal{G}_{s,\mathsf{M}}(\mathrm{I})\) 的诸拓扑的归纳极限拓扑。

设 \(M < M'\)，并设 B 是 \(\mathcal{G}_{s,M}(I)\) 中的（闭）单位球。我们将证明 B 是 Banach 空间 \(\mathcal{G}_{s,M'}(I)\) 的一个紧子集。显然 B 在 \(\mathcal{G}_{s,M'}(I)\) 中是闭的，故只需证明 B 在 \(\mathcal{G}_{s,M'}(I)\) 中是预紧的。设 \(\varepsilon > 0\)，并设 N 是使 \((M/M')^N \leqslant \varepsilon/2\) 的一个正整数。设 k 是一个正整数；当 f 取遍 B 时，诸函数 \(D^{k+1}f\) 所成之集在 \(\mathcal{C}(I)\) 中是有界的，因而当 f 取遍 B 时，诸函数 \(D^k f\) 所成之集在 \(\mathcal{C}(I)\) 中是相对紧的：这由有限增量定理（FVR, I, 第 23 页, 推论 1）以及 Ascoli 定理（GT, X, § 2, No.~5）得出。我们在 \(\mathcal{G}_{s,M}(I)\) 上借助下式定义一个范数

\[
q(f) = \sup_{\substack{0\leqslant n\leqslant \mathbf{N}\\ x\in \mathrm{I}}}\big|\mathrm{D}^{n}f(x)\big| / \mathrm{M}^{\prime n}(n!)^{s}  .
\]

上述论证表明，B 对于范数 q 相应的拓扑而言是预紧的；换言之，存在 B 的一个有限子集 C，使得对每个 \(f \in B\)，存在 \(g \in C\) 满足 \(q(f - g) \leqslant \varepsilon\)。最后，对每个 n \textgreater{} N，我们有

\[
\left| \mathrm{D} ^ {n} f (x) - \mathrm{D} ^ {n} g (x) \right| / \mathrm{M} ^ {\prime n} (n!) ^ {s} \leqslant 2 (\mathrm{M} / \mathrm{M} ^ {\prime}) ^ {n} \leqslant \varepsilon ,
\]

由此得到 \(\| f - g\| \leqslant \varepsilon\)。这就证明了 B 在 \(\mathcal{G}_{s,M'}(I)\) 中是预紧的。

空间 \(\mathcal{G}_s(\mathrm{I})\) 是诸空间 \(\mathcal{G}_{s,k}(\mathrm{I})\)（当 \(k\) 取遍 \(\mathbf{N}\) 时）的归纳极限；由命题 7（III, 第 6 页），\(\mathcal{G}_s(\mathrm{I})\) 的每个有界子集都含于诸空间 \(\mathcal{G}_{s,k}(\mathrm{I})\) 之一，并且在这个空间中是相对紧的。

* \(d\) )（全纯函数空间。）设 \(n \geqslant 1\) 是一个整数。对 \(\mathbf{C}^n\) 的每个开子集 U，\(\mathcal{H}(\mathrm{U})\) 表示在 U 中全纯的函数所成的空间，并赋予 U 中紧收敛拓扑。对 \(\mathbf{C}^n\) 的每个紧子集 L，\(\mathcal{H}(\mathrm{L})\) 表示在 L 的一个邻域中全纯的函数的芽所成的空间；我们给这个空间赋予使诸典范映射 \(\pi_{\mathrm{U}}:\mathcal{H}(\mathrm{U})\to\mathcal{H}(\mathrm{L})\) 连续的最细局部凸拓扑，其中 U 取遍 L 的开邻域所成之集。

对每个整数 \(m \geqslant 1\)，设 \(U_{m}\) 是 \(C^{n}\) 中与 L 的距离 \textless{} 1/m 的点所成之集。可以证明，典范映射 \(\pi_{U_{m}}\)（从 \(\mathcal{H}(U_{m})\) 到 \(\mathcal{H}(L)\) 中）是单射，并且从 \(\mathcal{H}(U_{m})\) 到 \(\mathcal{H}(U_{p})\) 中的限制映射对 \(p \geqslant m\) 是紧的。于是可以应用命题 7（III, 第 6 页）。设 A 是 \(\mathcal{H}(L)\) 的一个有界子集；则存在一个整数 \(m \geqslant 1\)，使得 A 由 L 的一个邻域中函数的芽组成，这些函数属于 \(\mathcal{H}(U_{m})\) 中的一个有界集 B。此外，映射 \(\phi\)（从 \(\mathcal{H}(L)\) 到一个拓扑空间 T 中）连续，当且仅当映射 \(\phi \circ \pi_{U}\)（从 \(\mathcal{H}(U)\) 到 T 中）对 L 的每个开邻域 U 都连续。*

